\documentclass[11pt]{article}
\usepackage{amsmath,amssymb,amsthm}
\usepackage[utf8]{inputenc}
\newcommand{\dd}{\mathrm{d}}
\newcommand{\ii}{\mathrm{i}}
\newcommand{\ee}{\mathrm{e}}
\newcommand{\Order}{\mathcal{O}}

\newtheorem{proposition}{Proposition}

\title{Direction D: Noise Effects on Higher-Order Kuramoto Synchronization}
\author{HaibiMathAgent}
\date{March 2026}

\begin{document}
\maketitle

\section{Noisy Higher-Order Kuramoto Model}

Adding Gaussian white noise:
\begin{equation}
  \dd\theta_i = \Bigl[\omega_i + \frac{K_2}{N}\sum_j\sin(\theta_j-\theta_i)
  + \frac{K_3}{N^2}\sum_{j,k}\sin(2\theta_j-\theta_k-\theta_i)\Bigr]\dd t
  + \sqrt{2\eta}\,\dd W_i,
\end{equation}
where $\eta$ is the noise intensity and $W_i$ are independent Wiener processes.

\section{Fokker--Planck Equation}

The single-oscillator density $f(\theta,\omega,t)$ satisfies:
\begin{equation}
  \frac{\partial f}{\partial t}
  + \frac{\partial}{\partial\theta}\bigl(f\,v\bigr)
  = \eta\,\frac{\partial^2 f}{\partial\theta^2},
\end{equation}
where $v[\theta,\omega,t]$ is the velocity field (same as deterministic case).

\section{Modified OA Reduction}

The Fourier expansion remains $f = \frac{g(\omega)}{2\pi}[1 + \sum_n
\hat{f}_n \ee^{\ii n\theta} + \text{c.c.}]$.

The noise adds a diffusion term: the $n$-th mode decays as $\ee^{-n^2\eta t}$.
On the OA manifold $\hat{f}_n = \alpha^n$, the modified evolution is:
\begin{equation}
  \dot{\alpha} = -(\ii\omega + \eta)\alpha
  + \frac{1}{2}\bigl[H - \bar{H}\alpha^2\bigr],
\end{equation}
where $H = K_2 Z_1 + K_3 Z_2\bar{Z}_1$ as before.

\paragraph{Key insight:} Noise simply shifts $\Delta \to \Delta + \eta$
in the Lorentzian case.  The radial dynamics become:
\begin{equation}
  \boxed{
  \dot{r} = -(\Delta+\eta)\,r + \frac{r(1-r^2)}{2}(K_2 + K_3 r^2).
  }
\end{equation}

\section{Consequences}

\begin{proposition}[Noise shifts the phase diagram uniformly]
  All critical manifolds shift by replacing $\Delta\to\Delta+\eta$:
  \begin{align}
    K_2^{\mathrm{c}}(\eta) &= 2(\Delta+\eta), \\
    \text{Explosive boundary:}\quad K_3 &= K_2\quad\text{(unchanged)}, \\
    \text{Fixed point eq:}\quad \Delta+\eta &= \frac{(1-r^2)}{2}(K_2+K_3 r^2).
  \end{align}
  The onset is shifted (higher $K_2$ needed), but the explosive boundary
  $K_3=K_2$ and the qualitative structure are noise-independent.
\end{proposition}

\paragraph{Physical interpretation:} Noise acts as an ``effective frequency
spread,'' broadening the distribution and making synchronization harder.  But
it does not change \emph{which type} of transition occurs---it only shifts
the critical point.

\section{Non-Gaussian (L\'{e}vy) Noise}

For L\'{e}vy noise with stability index $\alpha_L\in(0,2)$, the diffusion
term becomes a fractional operator:
\begin{equation}
  \eta\frac{\partial^2 f}{\partial\theta^2}
  \to \eta_L\frac{\partial^{\alpha_L} f}{\partial|\theta|^{\alpha_L}}.
\end{equation}

The $n$-th Fourier mode now decays as $\ee^{-|n|^{\alpha_L}\eta_L t}$
instead of $\ee^{-n^2\eta t}$.  On the OA manifold, this modifies the
effective damping to be mode-dependent, breaking the simple
$\Delta\to\Delta+\eta$ shift.

For the first mode ($n=1$, relevant for $r_1$):
\begin{equation}
  \dot{r} = -(\Delta+\eta_L)\,r + \frac{r(1-r^2)}{2}(K_2+K_3 r^2).
\end{equation}
This has the \emph{same form} as Gaussian noise for $r_1$!  The difference
appears only in higher-order parameters ($r_2$, $r_3$, etc.):
\begin{equation}
  Z_n \text{ decays as } \ee^{-|n|^{\alpha_L}\eta_L t},
\end{equation}
so $r_2$ decays faster than $r_1$ for $\alpha_L < 2$ (heavy-tailed noise
preferentially disrupts higher harmonics).

This provides a theoretical explanation for the Lee et al.\ (2025) finding
that L\'{e}vy noise produces ``extreme synchronization transitions'': the
$r_2 = r_1^2$ relation is broken more strongly, enabling novel collective
states.

\end{document}
